\documentclass[11pt]{article}
\usepackage[left=0.85in,right=0.85in,top=0.7in,bottom=0.7in]{geometry}
\usepackage[T1]{fontenc}
\usepackage{lmodern}
\usepackage{amsmath,amssymb,amsthm,mathtools,booktabs}
\usepackage[colorlinks=true,linkcolor=blue,citecolor=blue,urlcolor=blue]{hyperref}
\usepackage{microtype}
\newtheorem{theorem}{Theorem}
\newtheorem{proposition}[theorem]{Proposition}
\newtheorem{lemma}[theorem]{Lemma}
\newcommand{\Q}{\mathbb Q}
\newcommand{\Z}{\mathbb Z}
\newcommand{\R}{\mathbb R}
\newcommand{\dd}{\,\mathrm{d}}
\newcommand{\HH}[2]{H_{#1}^{(#2)}}
\title{Arithmetic saturation and toric Mellin models\\for odd-zeta approximation}
\author{}
\date{27 October 2026}
\begin{document}
\maketitle
\begin{abstract}
We give exact cofactor formulas for the arithmetic content of a tail-moment
construction and construct an algebraic two-parameter Mellin family whose
periods lie in $\Q+\Q\zeta(s)$ for every integer $s\ge2$. Partial fractions,
an explicit remainder estimate, and a pole-cancellation identity identify
these periods without a conjectural regulator comparison. The corresponding
integral determinant polynomials and several alternative normalizations are
computed exactly. The tests do not establish the required arithmetic
smallness. This note does not prove irrationality of every odd zeta value.
The accompanying Lean sources verify the stated summation and algebraic
identities and a conditional polynomial irrationality criterion.
\end{abstract}

\section{The arithmetic estimate under investigation}
For $s=2r+1>1$, let
\[
 \Delta_{N,a}=\det[\zeta(s)-\HH{a+i+j}{s}]_{0\le i,j<N}
             =A_{N,a}\zeta(s)+B_{N,a},\qquad A_{N,a}>0 .
\]
Set $M=a+2N-2$, $d_M=\operatorname{lcm}(1,\ldots,M)$ with $d_0=1$, $Q=d_M^s$,
$U=Q^NA_{N,a}$, $V=Q^NB_{N,a}$, and $g=\gcd(|U|,|V|)$.
The integer linear form is
\[
 \ell_{N,a}=\frac Ug\zeta(s)+\frac Vg
           =\frac Ug\,\frac{\Delta_{N,a}}{A_{N,a}}.
\]
The endpoint calculation in Appendix~\ref{app:endpoint} gives, for $N>r$,
\[
 \frac{(2r-1)!}{(N)_{2r}(N+a)_{2r}}
 \le \frac{\Delta_{N,a}}{A_{N,a}}
 \le \frac{(2r-1)!}{(N-r)_{2r}(N-r+a)_{2r}},
 \qquad (z)_k=\prod_{j=0}^{k-1}(z+j).
\]
It follows uniformly for arbitrary nonnegative shifts $a=a_N$ that
\begin{equation}
 \ell_{N,a_N}\longrightarrow0
 \quad\Longleftrightarrow\quad
 \frac{U/g}{N^{2r}(N+a_N)^{2r}}\longrightarrow0 .
 \label{eq:height}
\end{equation}
This is the unproved arithmetic estimate. In particular, the already
normalized coefficient, rather than an unnormalized determinant, is the
quantity that must be bounded. Since $U/g$ is a positive integer, the
displayed lower bound also prevents positive exponential decay in $N$
for this family when the shifts grow at most polynomially. Polynomial
smallness could still suffice; that assertion is not ruled out.

\section{An exact cofactor presentation}
Let $Z\in M_N(\Z)$ and write $J$ for the all-ones matrix. Define $B$ by
$B_{0j}=Z_{0j}$ and $B_{ij}=Z_{ij}-Z_{0j}$ for $i>0$, and put
$c_j=\operatorname{adj}(B)_{j0}$.
\begin{proposition}
For every scalar $x$,
\begin{equation}
 \det(Z+xJ)=x\sum_jc_j+\sum_j Z_{0j}c_j.
 \label{eq:cofactor}
\end{equation}
Consequently, if
$U=\det(Z+QJ)-\det Z$, $V=\det Z$, and
$d=\gcd_j|c_j|$, then $d$ divides both $U$ and $V$.
\end{proposition}
\begin{proof}
Subtract the first row from the other rows of $Z+xJ$. Only the first
row depends on $x$. Expanding along it gives \eqref{eq:cofactor};
the cofactors do not depend on that row. Each term in both sums is
divisible by $d$.
\end{proof}
The $c_j$ are the signed maximal minors of the lower $(N-1)\times N$
matrix. They satisfy $Lc=0$, by expanding a determinant with a repeated
row. When $d>0$, $\widehat c=c/d$ is its primitive integer kernel vector.
Writing
\[
 S=\sum_j\widehat c_j,\qquad C=\sum_jZ_{0j}\widehat c_j,\qquad
 h=\gcd(|QS|,|C|),
\]
gives $g=dh$ and, whenever $g>0$, the exact refinements
\begin{equation}
 g=dh,\qquad \frac Ug=\frac{QS}{h}.
 \label{eq:kernelheight}
\end{equation}
Thus divisibility of the maximal minors supplies an actual common
factor. A successful estimate must also control the functional $S$
and the remaining factor $h$. The Smith or Fitting interpretation of
the minors does not alone give that quantitative bound.
Lean verifies \eqref{eq:cofactor}, the lower-row kernel identity,
the common-divisor statement, and the full content factorization in
\eqref{eq:kernelheight}.

\section{A toric family with only the target zeta period}
Fix $s\ge2$, put $T=\prod_{i=0}^{s-1}x_i$, and define polynomial functions
\[
 A=\prod_{i=0}^{s-1}x_i^{\,i},\qquad
 B=\prod_{i=0}^{s-1}x_i^{\,s-1-i}.
\]
They give a morphism defined over $\Z$ and satisfy
\begin{equation}
 AB=T^{s-1}.
 \label{eq:toric}
\end{equation}
On the open cube, use the positive measure
$\omega=\dd x_0\cdots\dd x_{s-1}/(1-T)$ and set
\[
 I_s(j,k)=\int_{(0,1)^s} A^jB^k\,\omega,\qquad j,k\ge0 .
\]
The measure is finite: Tonelli's theorem and the geometric series give
$\int\omega=\sum_{q\ge0}(q+1)^{-s}=\zeta(s)$.

\begin{theorem}
For every $j,k\ge0$, $I_s(j,k)\in\Q+\Q\zeta(s)$. More precisely,
\begin{equation}
 I_s(j,j)=\zeta(s)-\HH{(s-1)j}{s}.
 \label{eq:diag}
\end{equation}
If $j\ne k$, let $e_i=ij+(s-1-i)k$ and
$c_i=\prod_{h\ne i}(e_h-e_i)^{-1}$. Then
\begin{equation}
 I_s(j,k)=-\sum_{i=0}^{s-1}c_iH_{e_i}\in\Q .
 \label{eq:rational}
\end{equation}
\end{theorem}
\begin{proof}
Expanding the positive geometric series under the integral gives
\[
 I_s(j,k)=\sum_{q=0}^{\infty}
              \prod_{i=0}^{s-1}(q+1+e_i)^{-1}.
\]
For $j=k$, all exponents equal $(s-1)j$, proving \eqref{eq:diag}.
Otherwise the exponents are distinct. Lagrange interpolation gives
\[
 \frac1{\prod_i(z+e_i)}=\sum_i\frac{c_i}{z+e_i},
 \qquad \sum_i c_i=0 .
\]
Summing through $q=L-1$ yields
\begin{equation}
 \sum_{q=0}^{L-1}\frac1{\prod_i(q+1+e_i)}
 =-\sum_i c_iH_{e_i}
   +\sum_i c_i\bigl(H_{L+e_i}-H_L\bigr).
 \label{eq:remainder}
\end{equation}
The remainder has the explicit bound
\[
 \left|\sum_i c_i(H_{L+e_i}-H_L)\right|
 \le\frac{\sum_i|c_i|e_i}{L+1},
\]
because $0\le H_{L+e_i}-H_L\le e_i/(L+1)$. It tends to zero.
This proves \eqref{eq:rational}, including convergence.
\end{proof}
Lean verifies the partial fractions, zero residue sum,
\eqref{eq:remainder}, its bound, the limit, and the resulting
\texttt{HasSum} statements. It also verifies \eqref{eq:diag} as a series
identity for the real part of the actual Riemann zeta function.

\subsection{An integral denominator and a pole-cancellation identity}
For distinct exponents in $[0,M]$, $d_M^{s-1}c_i$ is integral: every
nonzero difference has absolute value at most $M$ and divides $d_M$.
Since $d_MH_{e_i}$ is integral, $d_M^s I_s(j,k)\in\Z$.
On the diagonal, the same factor clears the rational term of
\eqref{eq:diag}. This gives an unconditional entrywise denominator
bound; it is not asserted to be sharp.

The defining relation \eqref{eq:toric} gives
\[
 \frac{AB-1}{1-T}=-\sum_{\ell=0}^{s-2}T^\ell.
\]
It follows, either by this identity or by shifting the convergent series,
that
\begin{equation}
 I_s(j+1,k+1)-I_s(j,k)
 =-\sum_{q=1}^{s-1}\prod_{i=0}^{s-1}(q+e_i)^{-1}.
 \label{eq:contiguity}
\end{equation}
Thus differences in the common diagonal direction have a finite rational
evaluation. Lean verifies both pole cancellation and
\eqref{eq:contiguity}. No assumption about a motivic lift is used in
these summation identities.

\section{Larger spaces and their integer determinant polynomials}
Brown's multi-parameter Mellin framework \cite{Brown} motivates using
many monomials at once. For the rectangular space
\[
 \mathcal M_n=\operatorname{span}_{\Q}
             \{A^jB^k:0\le j,k\le n\},
\]
let $G_n(X)$ be the matrix of pairings, with $\zeta(s)$ replaced by $X$
using \eqref{eq:diag}--\eqref{eq:rational}.
The monomials are independent: the exponent vectors of $A,B$ are
linearly independent. Consequently $G_n(\zeta(s))$ is positive definite.

The coefficient of $X$ in a pairing is one exactly when
$(j-k)+(j'-k')=0$. For every $j,k>0$, replace $A^jB^k$ by
$A^jB^k-A^{j-1}B^{k-1}$. This basis change has determinant one.
Equation \eqref{eq:contiguity} makes every pairing involving one such
difference rational. The remaining boundary basis
$1,A,\ldots,A^n,B,\ldots,B^n$ has $2n+1$ elements. Its coefficient
matrix is the permutation pairing opposite values of $j-k$.
The lower rational block is a positive Gram matrix, so
\[
 \deg\det G_n(X)=2n+1,
 \qquad \dim\mathcal M_n=(n+1)^2.
\]
This constructs a larger period space, while retaining the two periods
$1,\zeta(s)$.

The audit also uses balanced spaces
\[
 f_{t,u}=\prod_{i=0}^{s-1}
          x_i^{K+t+(i-r)u},\qquad
 s=2r+1,\quad |u|\le n,\quad 0\le t\le L,\quad K\ge rn .
\]
Products have either all equal or all distinct exponents, according as
$u+u'=0$ or not. Their Gram determinant has degree $2n+1$ and size
$(2n+1)(L+1)$. Indeed, replace $f_{t,u}$ for $t>0$ by
$f_{t,u}-f_{t-1,u}=(T-1)f_{t-1,u}$. Pairing a difference cancels
the denominator $1-T$, so the lower block is rational and positive.
The coefficient matrix on the remaining $f_{0,u}$ is the permutation
pairing $u$ with $-u$, proving the degree assertion.
A monic shifted-Chebyshev basis in $T$ and Hadamard's
inequality give the valid coarse analytic bound
\[
 0<\det G_{n,L}(\zeta(s))
   \le \zeta(s)^{(2n+1)(L+1)}\,2^{-2(2n+1)L^2}.
\]
Indeed each fixed-$u$ base monomial is bounded by one on the cube, and
the monic degree-$t$ factor has supremum at most $2^{1-2t}$ for $t\ge1$.
This estimate must be combined with the actual arithmetic multiplier.
It is not by itself an irrationality estimate.

Here that multiplier can be bounded row by row. Put $m=2n+1$,
$N=m(L+1)$, and
\[
 M_{t,u}=2K+t+L+r(|u|+n),\qquad
 D=\prod_{u=-n}^{n}\prod_{t=0}^{L}d_{M_{t,u}}^{\,s}.
\]
Every exponent in a product involving row $(t,u)$ lies between zero
and $M_{t,u}$. The denominator argument above therefore makes
$\widetilde P(X)=D\det G_{n,L}(X)$ an integer polynomial.
Let $g$ be the gcd of all its coefficients and set $P=\widetilde P/g$.
The analytic bound becomes the explicit arithmetic inequality
\begin{equation}
 0<P(\zeta(s))\le \frac Dg\,2^{N-2mL^2},
 \label{eq:mellinbudget}
\end{equation}
where $\zeta(s)\le2$ for $s\ge2$ follows from the integral test.
Consequently a proved bound
\[
 \frac{\log(D/g)+(N-2mL^2)\log2}{m}\longrightarrow-\infty
\]
along a sequence of these spaces would suffice for the polynomial
criterion below. This states both the full denominator cost and the
precise common-factor saving required; no such saving is assumed.

\subsection{The precise polynomial criterion}
Clear the coefficient denominators of a determinant polynomial and
divide out their entire integer gcd, obtaining $P_n\in\Z[X]$.
A rational change of basis multiplies the determinant by a nonzero
rational square, so this primitive polynomial is unchanged up to sign.

\begin{proposition}
Let $\alpha\in\R$. If for every integer $q>0$ there exists $P\in\Z[X]$
such that
\[
 0<|P(\alpha)|,\qquad q^{\deg P}|P(\alpha)|<1,
\]
then $\alpha$ is irrational.
\end{proposition}
\begin{proof}
If $\alpha=p/d$ with $d\ne0$, then
$d^{\deg P}P(p/d)$ is an integer. Nonvanishing forces its absolute value
to be at least one. Take $q=|d|$.
\end{proof}
For example, $\log|P_n(\alpha)|/\deg P_n\to-\infty$ with nonzero
evaluations is sufficient. Lean verifies the integer gap and the
proposition. This condition has not been established for the families
above; the degree factor cannot be omitted.

\section{Exact checks and the remaining bound}
All numbers below are finite certificates. Fraction-free integer
elimination constructs $U,V$; exact gcd reduction gives $u=U/g$.
The reported integer $b$ satisfies the checked inequality
\[
 u_{N,0}\ge 2^b N^{2(s-1)} .
\]
\begin{center}
\begin{tabular}{rrrr}
\toprule
Weight $s$ & $N=30$ & $N=40$ & $N=60$\\
\midrule
3 & 3325 & 6030 & 13726\\
5 & 6566 & 11876 & 27500\\
7 & 9797 & 17816 & 41008\\
9 & 12917 & 23580 & 54333\\
\bottomrule
\end{tabular}
\end{center}
The cofactor calculation was independently checked at twelve smaller
indices by the certificate $Bc=(\det B)e_0$, together with
\eqref{eq:cofactor}. It includes all common factors, not just primes
in the lcm scale.

For the balanced and toric spaces, all $57$ tested primitive determinant
polynomials have values strictly greater than one. The tested weights
are $3,5,7$, with matrix sizes up to $21$ and several base shifts.
A further pure-target multiple-Mellin construction was checked at $32$
weight/index pairs, through index $40$; its tested primitive absolute
values also exceed one. These computations do not rule out later
smallness or other spaces.

For completeness, this last construction is the explicit integer polynomial
\[
 F_{s,n}(x)=\sum_{k=0}^{sn}c_kx^k,\qquad
 c_k=(-1)^{sn-k}\binom{sn}{k}\binom{n+k}{n}^{s}.
\]
It gives the pure-target linear form
\[
 \int_0^1 F_{s,n}(x)\frac{(-\log x)^{s-1}}{(s-1)!(1-x)}\dd x
 =a_{s,n}\zeta(s)+b_{s,n},\quad
 a_{s,n}=\frac{(sn)!}{(n!)^s},\quad
 b_{s,n}=-\sum_kc_k\HH{k}{s}.
\]
The formula for $a_{s,n}=F_{s,n}(1)$ follows by taking the $sn$th
finite difference of the degree-$sn$ polynomial $\binom{n+k}{n}^{s}$,
whose leading coefficient is $(n!)^{-s}$. Its actual primitive
normalization is used in the audit; no decay is assumed.

\subsection{Rational certification of the evaluations}
The evaluations use rational intervals, not floating-point cancellation.
For integer $s\ge2$, Euler--Maclaurin summation through the $K$th
Bernoulli term gives
\[
 \zeta(s)=H_{M-1}^{(s)}+\frac{M^{1-s}}{s-1}+\frac1{2M^s}
 +\sum_{k=1}^K
    \frac{B_{2k}(s)_{2k-1}}{(2k)!\,M^{s+2k-1}}+R_K .
\]
The integral remainder and the Fourier bound
$|\widetilde B_{2K}(x)|\le |B_{2K}|$ imply
\[
 |R_K|\le
 \frac{|B_{2K}|(s)_{2K-1}}{(2K)!\,M^{s+2K-1}}.
\]
The bound is the absolute value of the final included term.
The scripts use $M=10K$, exact Bernoulli numbers, and rational interval
Horner evaluation. Increasing $K$ resolves the sign of any nonzero evaluation.
The derivation and the separate coarse integral enclosure provide
explicit checks on the normalization.

\subsection{What the files establish}
The Lean cofactor module verifies the expansion, kernel identity,
divisibility, and exact content factorization. The Mellin module verifies
the residue algebra, finite remainder, convergence, real zeta
identification, and toric contiguity. The polynomial module verifies
the conditional integer-polynomial criterion.
The integral geometry, positivity, degree argument, and Chebyshev
estimate are paper proofs here.
Neither \eqref{eq:height} nor the sufficient smallness condition for
the alternative polynomial families has been proved.
A full irrationality proof still requires one of these quantitative
arithmetic conclusions.

\paragraph{Scope of the arithmetic literature.}
Ferguson's positive-functional construction \cite{Ferguson} yields rational
Hankel approximations and explicitly stops short of arithmetic conclusions.
The denominator results of Krattenthaler--Rivoal \cite{KR} concern
specified hypergeometric coefficients. Applying their extra common
factors here would require an identity with those coefficients.
Brown's arithmetic refinements \cite{Brown} likewise require proved
denominator or congruence bounds for the chosen Mellin family. None of
these statements supplies the still-unproved estimate \eqref{eq:height}.

\appendix
\section{Derivation of the uniform endpoint estimate}
\label{app:endpoint}
The gamma integral and a nonnegative geometric series give
\[
 \zeta(s)-\HH{k}{s}=\int_0^1 x^k w_s(x)\dd x,
 \qquad w_s(x)=\frac{(-\log x)^{s-1}}{(s-1)!(1-x)}.
\]
These are moments of a finite positive measure. Every finite monomial
Gram matrix is positive definite, since a nonzero polynomial is nonzero
on an interior interval. In the basis
$1,(1-x),x(1-x),\ldots,x^{N-2}(1-x)$, the change-of-basis determinant
is $\pm1$. Only the first basis element is nonzero at $1$; the lower
block has entries
$\mu_k-2\mu_{k+1}+\mu_{k+2}=(k+1)^{-s}-(k+2)^{-s}$.
It is a positive Gram matrix and its determinant is $A_{N,a}>0$.
The Schur complement, or completion of the square in that basis, gives
\[
 E_{N,a}:=\frac{\Delta_{N,a}}{A_{N,a}}
 =\min_{\deg P<N,\ P(1)=1}\int_0^1P(x)^2x^aw_s(x)\dd x.
\]
For the Jacobi measure $x^b(1-x)^m\dd x$, with integers $m,b\ge0$,
the corresponding minimum is
\[
 \lambda_N^{(m,b)}(1)=\frac{m!(m+1)!}{(N)_{m+1}(N+b)_{m+1}}.
\]
Here is a direct check of the constant. The polynomials
\[
 J_j(x)=\frac{(-1)^j}{j!x^b(1-x)^m}
 \frac{\dd^j}{\dd x^j}[x^{b+j}(1-x)^{m+j}]
\]
are orthogonal by $j$ integrations by parts. Their squared norms and
endpoint values, obtained by the beta integral, are
\[
 h_j=\frac{(j+m)!(j+b)!}{(2j+m+b+1)j!(j+m+b)!},\qquad
 J_j(1)=\binom{j+m}{j}.
\]
All integration boundary terms vanish by the endpoint zero orders.
Cauchy--Schwarz gives the minimum
$1/\sum_{j<N}J_j(1)^2/h_j$, attained by coefficients proportional to
$J_j(1)/h_j$. The sum equals
$(N)_{m+1}(N+b)_{m+1}/[m!(m+1)!]$; subtracting its values at
$N=j+1$ and $N=j$ checks each summand.

For $s=2r+1$, put $q(x)=x^a(1-x)^{2r-1}/(2r)!$ and
$\rho=x^aw_s$. The elementary inequalities
$1-x\le-\log x\le(1-x)/x$ give $q\le\rho$ and
$x^{2r}\rho\le q$. The first yields the lower endpoint bound by
the Jacobi formula. For the upper bound restrict the minimum to
$P=x^rQ$, with $\deg Q<N-r$ and $Q(1)=1$, and use the second inequality.
Finally the scaled quantity $N^{2r}(N+a)^{2r}E_{N,a}$ lies between
\[
 (2r-1)!\left(\frac N{N+2r-1}\right)^{4r}
 \quad\hbox{and}\quad
 (2r-1)!\left(\frac N{N-r}\right)^{4r}.
\]
These bounds are uniform in $a\ge0$ and both tend to $(2r-1)!$,
proving \eqref{eq:height}.

\begin{thebibliography}{9}
\small
\bibitem{Brown}
F. Brown, \emph{Mellin transforms, transfinite diameter and rational
approximations of integrals}, 2026.
\url{https://arxiv.org/abs/2604.20741}.
\bibitem{Ferguson}
T. Ferguson, \emph{Rational Approximations via Hankel Determinants}, 2020.
\url{https://arxiv.org/abs/2003.10616}.
\bibitem{KR}
C. Krattenthaler and T. Rivoal,
\emph{Hyperg\'eom\'etrie et fonction z\^eta de Riemann},
Memoirs of the American Mathematical Society 186 (2007), no. 875.
\url{https://arxiv.org/abs/math/0311114}.
\end{thebibliography}
\end{document}
